%% ma: L11-B27-0016
%% lop: 11
%% chuong: 7
%% bai: 27
%% dang: 11.27.12
%% loai: TLN
%% muc: VD
%% dap_an: "9"
%% nguon: "(NBV)-ĐỀ SỐ 14-MỖI NGÀY 1 ĐỀ THI 2025.pdf (D14, MỖI NGÀY 1 ĐỀ - Đề số 14), Phần III Câu 2"
%% kien_thuc: "Bài 25 (góc nhị diện), Bài 23 (đường thẳng vuông góc với mặt phẳng), Bài 27 (thể tích khối chóp)"
%% trang_thai: chinh-thuc
%% ghi_chu: "Đã duyệt 15/10/2026. Đáp án gốc đúng (đã kiểm bằng tọa độ). Đề gốc không có hình toán nên không vẽ hình; ảnh minh họa bên cạnh chỉ là trang trí."
\begin{ex}
Cho hình chóp $S.ABCD$ có đáy $ABCD$ là hình vuông cạnh $3$, $SA$ vuông góc với mặt phẳng đáy. Biết số đo của góc nhị diện $[B,SC,D]$ bằng $120^\circ$. Tính thể tích khối chóp $S.ABCD$.
\shortans{9}
\loigiai{Đặt $SA=h$. Vì $BC\perp(SAB)$ nên $BC\perp SB$; tương tự $CD\perp SD$. Hai tam giác vuông $SBC$ và $SDC$ bằng nhau (cạnh huyền $SC$ chung, $BC=DC$), nên kẻ $BE\perp SC$ thì $DE\perp SC$ và $\widehat{BED}=120^\circ$ là góc phẳng của góc nhị diện.
Gọi $O=AC\cap BD$; tam giác $EBD$ cân tại $E$ nên $\widehat{BEO}=60^\circ$ và $BE=\dfrac{BO}{\sin60^\circ}=\dfrac{3\sqrt2/2}{\sqrt3/2}=\sqrt6$.
Trong tam giác $SBC$ vuông tại $B$: $BE=\dfrac{SB\cdot BC}{SC}=\dfrac{3\sqrt{h^2+9}}{\sqrt{h^2+18}}=\sqrt6\Rightarrow 9(h^2+9)=6(h^2+18)\Rightarrow h=3$.
Vậy $V=\dfrac13\cdot3^2\cdot3=9$.}
\end{ex}
